% ══════════════════════════════════════════════════════════════════════════════
\chapter*{Conclusion}
\addcontentsline{toc}{chapter}{Conclusion}
% ══════════════════════════════════════════════════════════════════════════════

Ce projet a permis de construire, de la donnée brute au moteur opérationnel, un système
complet d'indexation et de recherche d'information sur le corpus ADIT. Les six étapes
correspondent aux briques fondamentales de toute architecture RI : constitution du
corpus, contrôle du vocabulaire, réduction morphologique, tolérance aux fautes, analyse
de requête, et évaluation. Le moteur final atteint P\,=\,R\,=\,F1\,=\,1,000 sur le jeu
de 10 requêtes, avec un temps de réponse moyen de 0,63\,ms, et une couverture de tests
de 93\,\% sur 369 tests unitaires.

D'un point de vue méthodologique, deux enseignements principaux ressortent. Premièrement,
l'évaluation honnête d'un système RI exige un ground truth \emph{indépendant} du système :
c'est cette indépendance qui a permis d'identifier deux bugs réels (normalisation des
rubriques, verbe de cadrage parasite) qui auraient été invisibles avec un oracle interne.
Deuxièmement, les petits modèles de langue (ici \texttt{fr\_core\_news\_sm}) sont
souvent insuffisants pour des domaines spécialisés : 0/12 entités DATE reconnues sur
les requêtes de test ont conduit à préférer les expressions régulières hiérarchisées,
plus reproductibles et plus rapides sur ce corpus.

Les pistes d'amélioration identifiées sont l'indexation par trigrammes pour le correcteur
(robustesse aux erreurs en début de mot), le passage à BM25 pour le scoring (normalisation
par longueur de document), et l'utilisation d'un modèle de langue de type CamemBERT pour
une analyse sémantique plus fine des requêtes complexes.
